\documentclass[11pt,a4paper]{article}
\usepackage[margin=25mm]{geometry}
\usepackage{amsmath,amssymb,amsthm,booktabs,hyperref}
\setlength{\emergencystretch}{2em}
\newtheorem{theorem}{Theorem}
\newtheorem{lemma}[theorem]{Lemma}
\newcommand{\gen}[1]{\langle#1\rangle}
\title{Disconnected strings in two Coxeter families}
\author{Human Compression project: F006 recovery}
\date{}
\begin{document}\maketitle
A string C-marking is an ordered list of nonidentity involutions generating a
finite group, with nonadjacent entries commuting and with
$G_I\cap G_J=G_{I\cap J}$ for every two sets of positions, where $G_I$
is generated by the entries in $I$. We allow adjacent entries to commute.
Such a marking is called disconnected when at least one adjacent product has
order two. These conventions are the literal ones in Chen and Rowley
\cite{source}; their discussion of $B_n$ expressly includes disconnected strings.

Their Theorem~3(i) and Table~1 require a qualification under this convention.
The two cases addressed here are $I_2(m)$ and $H_3$. The eight connected
$H_3$ markings in the source table are recovered, but there are also three
rank-four classes. This note does not certify the other exceptional groups or
the $D_n$ results. No new open-conjecture answer is asserted.

\section{The dihedral group}
\begin{theorem}
Let $m\ge3$ and let $D=\gen{r,s\mid r^m=s^2=1,\ srs=r^{-1}}$.
With disconnected strings allowed,
\[
 r_{\max}(D)=
 \begin{cases}3,&m\equiv2\pmod4,\\2,&\text{otherwise.}\end{cases}
\]
Every rank-three string C-marking of $D$ is disconnected.
\end{theorem}
\begin{proof}
The involutions are the reflections $r^i s$ and, when $m$ is even, the
central involution $z=r^{m/2}$. Two distinct reflections commute precisely
when their product is $z$. The intersection property makes every generator
independent of the others, and in particular makes all generators distinct.

There is no marking of rank at least four. If $z$ occurs, at least three
other entries are reflections; two of their positions are nonadjacent.
They commute, and their product is $z$, contrary to independence. If $z$
does not occur, the first reflection commutes with both the third and the
fourth. A reflection has at most one distinct commuting reflection, again
a contradiction.

Suppose a rank-three marking consists entirely of reflections. Its end
entries commute, so $m$ is even. Change the name of the first reflection to
$s$; the last is $zs$, and the middle has the form $r^t s$.
Generation requires $\gcd(t,m/2)=1$. If $m/2$ is even, then $t$ is odd,
so $\gcd(t,m)=1$ and the first two entries already generate $D$.
If $m/2$ is odd, either $t$ is odd, with the same consequence, or
$m/2-t$ is odd and coprime to $m$, in which case the last two entries
generate $D$. Every case contradicts independence.

A rank-three marking must therefore consist of $z$ and two reflections
$a,b$. Put $L=\gen{a,b}$. Independence says $z\notin L$. If
$h=|ab|$, the rotation subgroup of $L$ has order $h$ and contains $z$
exactly when $h$ is even. Thus $h$ is odd, and generation of $D$ gives
$m=2h$. The central entry cannot be in the middle, since the two end
reflections would then commute and generate it. The marking is consequently
disconnected.

Conversely, if $m=2h$ with odd $h\ge3$, then
$D=\gen{z}\times\gen{s,r^2s}$. The second factor is the dihedral group of
order $2h$. Its two distinct generating reflections have the rank-two
intersection property. For the marking $(z,s,r^2s)$, every subgroup generated
by a subset is the direct product of its two coordinate subgroups, so the
intersection property follows coordinate by coordinate. Its type is
$\{2,h\}$. The usual pair $(s,rs)$ gives rank two in every case.
\end{proof}

For example $I_2(6)$ has the disconnected rank-three marking of type
$\{2,3\}$. The rank-two assertion is valid when connected strings alone
are intended.

\section{The icosahedral group}
We retain the classical finite Coxeter order $|W(H_3)|=120$. Its
identification with the following permutation group is verified explicitly.
Set
\[
 a=(1\ 2)(3\ 4),\quad b=(1\ 2)(3\ 5),\quad
 c=(1\ 3)(2\ 4),\quad z=(6\ 7).
\]
The products $ab,bc,ac$ have orders $3,5,2$, respectively. Hence
$\gen{a,b,c}\le A_5$ has subgroups of orders $6,10,4$; its order is
divisible by their least common multiple, sixty, and it is $A_5$.
The product $abc$ has order five. Therefore the Coxeter triple
$(za,zb,zc)$, of type $\{3,5\}$, generates $A_5\times\gen z$: its
product has fifth power $z$. The resulting surjection from $W(H_3)$ is
an isomorphism by the retained order. No isomorphism with a group of the
same order alone is presumed.

\begin{lemma}
The marking $(z,a,b,c)$ is a rank-four string C-marking of $W(H_3)$,
of type $\{2,3,5\}$.
\end{lemma}
\begin{proof}
The three rank-two subgroups for $(a,b,c)$ have orders $6,10,4$.
Their pairwise intersections have orders dividing two and contain the
respective common generator. Thus each intersection is precisely the
common subgroup. Together with their cyclic subgroups and the trivial
subgroup this checks the intersection property for all eight subsets.
The end involutions $a,c$ commute. Finally adjoining $z$ as a separate
end factor preserves every intersection coordinate by coordinate, as in
the preceding theorem.
\end{proof}

\begin{theorem}[Complete small census]
Up to group automorphisms and reversal, the string C-markings of $W(H_3)$
have the following counts. Parentheses give the number of self-dual classes.
\[
\begin{array}{c|ccc}
 &\text{rank }3&\text{rank }4&\text{rank }\ge5\\\hline
\text{connected}&8\ (1)&0&0\\
\text{disconnected}&0&3\ (0)&0
\end{array}
\]
In particular its maximum rank is four with both conventions included,
and three when connected strings alone are considered.
\end{theorem}
\begin{proof}[Computational proof with the complete finite search specified]
Use the 60 even permutations of $\{1,\ldots,5\}$, each independently
fixing or exchanging $6,7$, to obtain 120 elements. Multiplication is ordinary
permutation composition. There are 31 nonidentity involutions.
The program \texttt{check\_h3.py} constructs this whole multiplication table;
it uses no group classification or canonical-label library.

It grows every ordered prefix of involutions. When appending $x$, it requires
$x$ to lie outside the subgroup already generated, checks its commutation
with all earlier nonadjacent entries, and tests the intersection property
for every pair of subsets of the enlarged prefix. Subgroups are computed
by repeatedly multiplying known elements by the specified generators until
no new element remains. These are exact finite set operations.
Every prefix of a string C-marking has the intersection property and
independent generators, so no possible marking is lost by these rejections.
When the generated subgroup is the whole group, the marking is recorded;
it cannot be extended independently. Each proper extension at least doubles
the subgroup order. Thus any marking has $2^r\le120$, and it suffices to
search through rank six. The numbers of valid prefixes reached at ranks
one through six are
\[
                  31,\quad930,\quad3720,\quad720,\quad0,\quad0.
\]

For completeness of the equivalence calculation, every automorphism of
$A_5\times C_2$ preserves its derived subgroup $A_5$ and its centre $C_2$.
Here $A_5$ is perfect because a commutator of two double transpositions is
a three-cycle, and conjugate three-cycles generate $A_5$; its centre is
trivial. Hence the automorphism group is $\operatorname{Aut}(A_5)$.
There are five Sylow two subgroups of $A_5$, one Klein four subgroup fixing
each point. An automorphism permutes these five subgroups faithfully:
if it fixes them all, conjugation by $g$ and by its image have the same
permutation on them for every $g$, and the natural $A_5$ action is faithful.
Thus $\operatorname{Aut}(A_5)$ embeds in $S_5$. Conversely conjugation
by $S_5$ supplies 120 distinct automorphisms. These are all of them.

The program therefore takes the least tuple among these 120 conjugates
and their reversals. A class is self-dual exactly when its reversal is
among its 120 conjugates without reversal. It tests connectedness by
whether all adjacent product orders exceed two. The full representative
lists and the complete execution give the displayed counts. In particular,
no truncated search or sampled list supplies the zero columns. Correct
execution of the supplied Python integer and permutation program remains
the explicit computational premise; this is not a Lean formalization.
\end{proof}

\paragraph{Attribution and scope.}
The source's definition, rank claims and eight connected $H_3$ classes are
those of Chen and Rowley. The dihedral structure, direct-product intersection
argument, classical order of $H_3$, and automorphism description of $A_5$
are standard inputs or elementary arguments given here. The present work
makes the disconnected boundary explicit and gives an independent complete
small enumeration. This is a correction to the literal scope, not a claim
that the source's connected census is wrong. The other five exceptional
census rows, the $D_n$ constructions and all whole-manuscript admission
obligations remain outstanding. No priority claim or conjecture credit is
attached to this bounded result.

\begin{thebibliography}{1}
\bibitem{source} M. H. W. Chen and P. Rowley,
\emph{Abstract Regular Polytopes of Finite Irreducible Coxeter Groups},
Electronic Journal of Combinatorics \textbf{32}(4) (2025), \#P4.18,
\href{https://doi.org/10.37236/13763}{doi:10.37236/13763},
fixed thirteen-page published edition, especially Theorem~3 and Table~1.
\end{thebibliography}
\end{document}
